%%%%%%%%%%%%%%%%%%%%%%%%%%%%%%%%%%%%%%%%%%%%%%%%%%%%%%%%%%%%%%%%%%%%%%
% Lớp 9 | Chương V. Đường tròn | Vị trí tương đối của đường thẳng và đường tròn
% Dạng 9H5?16-1: Xác định vị trí tương đối của đường thẳng và đường tròn
%%%%%%%%%%%%%%%%%%%%%%%%%%%%%%%%%%%%%%%%%%%%%%%%%%%%%%%%%%%%%%%%%%%%%%


% nguồn: 2026-11-21_lop9_chuong5_bai14_17 (nhập 21/11/2026)
\begin{ex}%[9H5N16-1]%[ma:f354ec49]
% nguon: SGK Toán 9 KNTT tập 1 – Bài 16, Vị trí tương đối của đường thẳng và đường tròn
	Cho đường tròn $(O; 5\text{ cm})$ và đường thẳng $a$ có khoảng cách từ tâm $O$ đến $a$ bằng $7\text{ cm}$. Vị trí tương đối của đường thẳng $a$ và đường tròn $(O)$ là
	\choice
	{Cắt nhau}
	{Tiếp xúc nhau}
	{\True Không giao nhau}
	{Đi qua tâm $O$}
	\loigiai{
		Vì khoảng cách từ tâm $O$ đến đường thẳng $a$ là $d = 7\text{ cm} > R = 5\text{ cm}$ nên đường thẳng $a$ và đường tròn $(O)$ không giao nhau.
	}
\end{ex}

% nguồn: 2026-11-21_lop9_chuong5_bai14_17 (nhập 21/11/2026)
\begin{ex}%[9H5H16-1]%[ma:2bdead0d]
% nguon: SGK Toán 9 KNTT tập 1 – Bài 16, Vị trí tương đối của đường thẳng và đường tròn
	Cho đường tròn $(O; 6\text{ cm})$ và ba đường thẳng $a, b, c$ có khoảng cách đến tâm $O$ lần lượt là $d_a = 4\text{ cm}$, $d_b = 6\text{ cm}$, $d_c = 8\text{ cm}$. Xét tính đúng, sai của các khẳng định sau:
	\choiceTF
	{\True Đường thẳng $a$ cắt đường tròn $(O)$ tại hai điểm phân biệt}
	{\True Đường thẳng $b$ là tiếp tuyến của đường tròn $(O)$}
	{Đường thẳng $c$ và đường tròn $(O)$ có đúng một điểm chung}
	{\True Độ dài dây cung tạo bởi đường thẳng $a$ và đường tròn $(O)$ bằng $4\sqrt{5}\text{ cm}$}
	\loigiai{
		\begin{itemchoice}
			\itemch Đúng. Vì $d_a = 4\text{ cm} < R = 6\text{ cm}$ nên đường thẳng $a$ cắt đường tròn $(O)$ tại hai điểm phân biệt.
			\itemch Đúng. Vì $d_b = 6\text{ cm} = R$ nên đường thẳng $b$ tiếp xúc với đường tròn $(O)$, tức là tiếp tuyến của $(O)$.
			\itemch Sai. Vì $d_c = 8\text{ cm} > R = 6\text{ cm}$ nên đường thẳng $c$ và đường tròn $(O)$ không có điểm chung nào.
			\itemch Đúng. Độ dài dây cung cắt bởi đường thẳng $a$ là $2\sqrt{R^2 - d_a^2} = 2\sqrt{6^2 - 4^2} = 2\sqrt{20} = 4\sqrt{5}\text{ cm}$.
		\end{itemchoice}
	}
\end{ex}

% nguồn: 2026-11-21_lop9_chuong5_bai14_17 (nhập 21/11/2026)
\begin{bt}%[9H5V16-1]%[ma:6c2dbd8b]
% nguon: SGK Toán 9 KNTT tập 1 – Bài 16, Vị trí tương đối của đường thẳng và đường tròn
	Trong mặt phẳng tọa độ $Oxy$, cho điểm $A(3; 4)$. Gọi $R$ là bán kính của đường tròn tâm $A$ sao cho đường tròn đó tiếp xúc với trục tung $Oy$. Giá trị của $R$ bằng bao nhiêu?
	\shortans{3}
	\loigiai{
		Khoảng cách từ điểm $A(3; 4)$ đến trục tung $Oy$ bằng $|x_A| = |3| = 3$.\\
		Để đường tròn $(A; R)$ tiếp xúc với trục tung $Oy$ thì $R = d(A, Oy) = 3$.
	}
\end{bt}

% nguồn: 2026-11-21_lop9_chuong5_bai16_17_bs (nhập 21/11/2026)
\begin{ex}%[9H5N16-1]%[ma:e42da1c9]
% nguon: SGK Toán 9 KNTT tập 1 – Bài 16, Vị trí tương đối của đường thẳng và đường tròn
	Nếu đường thẳng $d$ cắt đường tròn $(O; R)$ thì khoảng cách $h$ từ tâm $O$ đến đường thẳng $d$ thỏa mãn
	\choice
	{\True $h < R$}
	{$h = R$}
	{$h > R$}
	{$h \ge 2R$}
	\loigiai{
		Đường thẳng $d$ cắt đường tròn $(O; R)$ khi và chỉ khi khoảng cách từ tâm $O$ đến $d$ nhỏ hơn bán kính, tức là $h < R$.
	}
\end{ex}

% nguồn: 2026-11-21_lop9_chuong5_bai16_17_bs (nhập 21/11/2026)
\begin{ex}%[9H5N16-1]%[ma:b03052ea]
% nguon: SGK Toán 9 KNTT tập 1 – Bài 16, Vị trí tương đối của đường thẳng và đường tròn
	Cho đường tròn $(O; 8\text{ cm})$ và đường thẳng $a$ có khoảng cách từ tâm $O$ đến $a$ bằng $8\text{ cm}$. Số điểm chung của đường thẳng $a$ và đường tròn $(O)$ là
	\choice
	{$0$}
	{\True $1$}
	{$2$}
	{Vô số}
	\loigiai{
		Vì khoảng cách từ tâm $O$ đến đường thẳng $a$ bằng bán kính $R = 8\text{ cm}$ nên $a$ tiếp xúc với $(O)$, do đó chúng có đúng $1$ điểm chung.
	}
\end{ex}

% nguồn: 2026-11-21_lop9_chuong5_bai16_17_bs (nhập 21/11/2026)
\begin{ex}%[9H5H16-1]%[ma:369e301c]
% nguon: SGK Toán 9 KNTT tập 1 – Bài 16, Vị trí tương đối của đường thẳng và đường tròn
	Trong mặt phẳng tọa độ $Oxy$, cho đường tròn tâm $M(4; 3)$ bán kính $R = 4$. Xét tính đúng, sai của các khẳng định sau:
	\choiceTF
	{\True Khoảng cách từ tâm $M$ đến trục hoành $Ox$ bằng $3$}
	{\True Khoảng cách từ tâm $M$ đến trục tung $Oy$ bằng $4$}
	{Đường tròn $(M; 4)$ không giao nhau với trục hoành $Ox$}
	{\True Đường tròn $(M; 4)$ tiếp xúc với trục tung $Oy$ và cắt trục hoành $Ox$ theo một dây cung có độ dài bằng $2\sqrt{7}$}
	\loigiai{
		\begin{itemchoice}
			\itemch Đúng. Khoảng cách từ $M(4; 3)$ đến trục hoành $Ox$ là $|y_M| = 3$.
			\itemch Đúng. Khoảng cách từ $M(4; 3)$ đến trục tung $Oy$ là $|x_M| = 4$.
			\itemch Sai. Vì khoảng cách từ $M$ đến $Ox$ là $3 < R = 4$ nên đường tròn $(M; 4)$ cắt trục hoành $Ox$ tại hai điểm phân biệt.
			\itemch Đúng. Vì $d(M, Oy) = 4 = R$ nên $(M; 4)$ tiếp xúc với $Oy$; độ dài dây cung cắt trên $Ox$ là $2\sqrt{R^2 - d^2} = 2\sqrt{4^2 - 3^2} = 2\sqrt{7}$.
		\end{itemchoice}
	}
\end{ex}

% nguồn: 2026-11-21_lop9_chuong5_bai16_17_bs (nhập 21/11/2026)
\begin{ex}%[9H5H16-1]%[ma:7171ce65]
% nguon: SGK Toán 9 KNTT tập 1 – Bài 16, Vị trí tương đối của đường thẳng và đường tròn
	Cho đường tròn $(O; 5\text{ cm})$ và đường thẳng $d$ cách tâm $O$ một khoảng $h = 3\text{ cm}$. Xét tính đúng, sai của các khẳng định sau:
	\choiceTF
	{\True Đường thẳng $d$ cắt đường tròn $(O)$ tại hai điểm phân biệt $A, B$}
	{\True Độ dài dây cung $AB$ bằng $8\text{ cm}$}
	{Nếu dịch chuyển đường thẳng $d$ ra xa tâm $O$ thêm $1\text{ cm}$ thì đường thẳng mới tiếp xúc với $(O)$}
	{\True Diện tích tam giác $OAB$ bằng $12\text{ cm}^2$}
	\loigiai{
		\begin{itemchoice}
			\itemch Đúng. Vì $h = 3 < R = 5$ nên $d$ cắt $(O)$ tại hai điểm phân biệt $A, B$.
			\itemch Đúng. Độ dài dây $AB = 2\sqrt{5^2 - 3^2} = 2 \cdot 4 = 8\text{ cm}$.
			\itemch Sai. Khi dịch chuyển ra xa tâm thêm $1\text{ cm}$ thì khoảng cách mới là $3 + 1 = 4\text{ cm} < 5\text{ cm}$, đường thẳng vẫn cắt đường tròn.
			\itemch Đúng. Diện tích tam giác $OAB$ là $\dfrac{1}{2} \cdot AB \cdot h = \dfrac{1}{2} \cdot 8 \cdot 3 = 12\text{ cm}^2$.
		\end{itemchoice}
	}
\end{ex}

% nguồn: 2026-11-21_lop9_chuong5_bai16_17_bs (nhập 21/11/2026)
\begin{bt}%[9H5V16-1]%[ma:64600a4b]
% nguon: SGK Toán 9 KNTT tập 1 – Bài 16, Vị trí tương đối của đường thẳng và đường tròn
	Trong mặt phẳng tọa độ $Oxy$, cho điểm $B(-5; 6)$. Gọi $R$ là bán kính của đường tròn tâm $B$ sao cho đường tròn đó tiếp xúc với trục hoành $Ox$. Giá trị của $R$ bằng bao nhiêu?
	\shortans{6}
	\loigiai{
		Khoảng cách từ điểm $B(-5; 6)$ đến trục hoành $Ox$ bằng $|y_B| = |6| = 6$.\\
		Để đường tròn $(B; R)$ tiếp xúc với trục hoành $Ox$ thì $R = d(B, Ox) = 6$.
	}
\end{bt}

% nguồn: 2026-11-21_lop9_chuong5_bai16_17_bs (nhập 21/11/2026)
\begin{bt}%[9H5V16-1]%[ma:b585d602]
% nguon: SGK Toán 9 KNTT tập 1 – Bài 16, Vị trí tương đối của đường thẳng và đường tròn
	Cho đường tròn $(O; 10\text{ cm})$ và đường thẳng $d$ cách tâm $O$ một khoảng bằng $6\text{ cm}$. Đường thẳng $d$ cắt đường tròn $(O)$ tại hai điểm $A$ và $B$. Tính diện tích tam giác $OAB$ (đơn vị $\text{cm}^2$).
	\shortans{48}
	\loigiai{
		Gọi $H$ là hình chiếu của $O$ trên $d$, ta có $OH = 6\text{ cm}$ và $AH = \sqrt{OA^2 - OH^2} = \sqrt{10^2 - 6^2} = 8\text{ cm}$.\\
		Suy ra $AB = 2AH = 16\text{ cm}$.\\
		Diện tích tam giác $OAB$ là $S = \dfrac{1}{2} \cdot AB \cdot OH = \dfrac{1}{2} \cdot 16 \cdot 6 = 48\text{ cm}^2$.
	}
\end{bt}
